\chapter{Dead Ends: What Did Not Work, and Why}
\label{ch:sackgassen}

% Sources: verworfen.md X1-X9, counterexamples.md, lab notes.
% Tone: narrative, honest, one clear lesson per dead end.

Research reports like to tell only the hits. This chapter does the
opposite: it collects the roads we walked down and walked back up
again. Not out of contrition, but because the dead ends of this
project have revealed more about Horimetrik than many a theorem --
and because at least one of them later turned out to be a detour
rather than a wrong turn.

\section{The $n!/2$ debacle, or: three data points are not a pattern}

\begin{sackgasse}
\emph{Idea:} At $n=5,6,7$ the shell-jumper sequence showed the
values $60, 360, 2520$ -- exactly $n!/2$. A formula!
\emph{Obstacle:} The fourth data point: $21168\ne20160$.
\emph{Lesson:} Patterns built from a few terms are hypotheses, not
results.
\end{sackgasse}

This story has a remarkable epilogue: much later, an exact digit
formula for the same sequence fell out of the sky -- and it
explains in hindsight \emph{why} the false formula held for three
terms. For all three hits the digit formula collapses to a single
summand, and that summand equals exactly $n!/2$: for $120$ and
$5040$ because they have single-digit representations in their own
shells, for $720$ because the first maximal digit cuts the formula
off right after its first contributing summand. The coincidence was
not noise but a genuine structural feature -- only one of the
digits of $n!$, not of the sequence itself. This is how dead ends
sometimes rehabilitate themselves: yesterday's mistake became the
footnote of a theorem.

\section{The small-world trap, six times over}

The same trap caught us six times in the course of this project, in
ever new disguises: a bound on the compactification commutator
looked like $\pm1$ up to the seventh world (in truth it is
unbounded on both sides); its successor, a one-sided bound,
withstood random sampling up to the twenty-first world and fell
only to \emph{constructed} witnesses; an element of the
compactification monoid acted idempotent up to the fourth world,
and the smallest counterexample lives in the fifth; and the $n!/2$
pattern you already know. Cases five and six happened at the upper
side of the depth law within a single evening: first a pretty
triangular-number hypothesis died on its own candidate sweep,
before it was even written down -- then the threshold that this
sweep delivered turned out, in its turn, to be an artefact of too
shallow a search. The common lesson is more precise than the usual
``test more'':

\begin{center}
\emph{In the Hori space the interesting counterexamples grow
factorially slowly.\\ Whoever claims bounds must construct, not
sample.}
\end{center}

Incidentally, the same number haunts several of these stories: the
six -- target of the smallest shell jump (the five, extended by one
zero, becomes the six), the value-smallest separating witness of
the monoid, the start of the first interesting shell. It is the
house spirit of this book.

\section{Discarded application dreams}

\begin{sackgasse}
\emph{Idea:} Hori addresses instead of IPv6, Hori heaps instead of
binary heaps.
\emph{Obstacle:} Within a single horizon the Hori space is
structurally rigid -- address spaces need local elasticity, and the
heap pays for the growing branching factor with expensive child
comparisons.
\emph{Lesson:} The system's strength is the global factorial
ordering, not local flexibility. Applications must fit the
structure, not the other way round.
\end{sackgasse}

\begin{sackgasse}
\emph{Idea:} Cardinality arguments -- is the Hori space ``bigger''
than the natural numbers?
\emph{Obstacle:} It is countable; any enumeration flattens it.
\emph{Lesson:} The added value lies not in cardinality but in
structure per value. A chessboard does not become more interesting
when you number its squares -- and not more boring either.
\end{sackgasse}

The third application dream came from outside, was the most
ambitious -- and died on the very first theorem of this book.

\begin{sackgasse}
\emph{Idea:} Cryptography from capacity (an external proposal, July
2026). A secret shape is pushed through $k$ leading zeros and
evaluated value-compactly; only the bare value is published. The
attacker, so the hope went, founders on the difference overflow --
the way back via antidifferences blows the capacity -- and the side
fidelity theorem blocks every hybrid shortcut.
\emph{Obstacle:} The way back needs neither antidifferences nor
mixed arithmetic. From $y=P(N)$, repeated division with remainder
pulls the digits down one by one: $c_0=y\bmod N$ is \emph{exact}
(since $c_0\le r<N$), then $y_1=(y-c_0)/N$ delivers the next round
with divisor $N{-}1$, and so on -- a handful of divisions, and the
``secret'' shape lies in the open. This is no ingenious attack but
the bijection from Chapter~\ref{ch:horizonte} at a shifted
evaluation point: the weights $1,N,N(N{-}1),\ldots$ again form a
mixed-radix system, and the digit bounds stay obediently below the
divisors.
\emph{Lesson:} The factorial ordering does not scramble, it
\emph{sorts} -- the very bijectivity that makes the system so clean
forbids it cryptographic hardness. Two clarifications for the
record: the side fidelity theorem forbids mixed \emph{semirings},
not mixed \emph{algorithms} -- an attacker may divide, tabulate,
and switch representations without ever mixing an arithmetic. And a
one-way function needs a secret; Horimetrik has none, it is
deliberately bright everywhere. What remains is the more modest
role: the horizon calculus of
Chapter~\ref{ch:strukturpolynome} works as a \emph{measuring
instrument} for capacity and noise growth in other people's
constructions -- not as a source of hardness of its own.
\end{sackgasse}

\section{The tropical mirage -- and its resolution}

The most instructive dead end deserves the finale. Early on it was
noticed that under addition the horizons \emph{want} to behave like
a maximum, and under multiplication like a sum -- the behaviour of
so-called tropical algebra. The direct attempt to compute that way
failed on an overflow problem: the value sum can burst the maximum
horizon. Discarded.

Months later -- in this project's reckoning of time: hours -- the
idea returned through the back door. The horizon splits into a
mandatory part and a freely chosen surplus, and \emph{on the
surplus} the tropical arithmetic works flawlessly: the mandatory
part absorbs all overflows, the memory computes with maximum and
sum. The idea was never wrong. It merely stood in the wrong
coordinate system.

\begin{oberflaeche}
If this chapter has one message, it is this: discarding is not
defeat, it is bookkeeping. Of our nine documented dead ends (the
working register lists them as X1 through X9), one later became a
theorem -- the tropical one --, one became the footnote of a
theorem -- the $n!/2$ debacle --, and the rest became precise
warning signs. None of that would have happened had we deleted our
failed attempts instead of dating them.
\end{oberflaeche}
